\documentclass[11pt]{article}

\usepackage[margin=1in]{geometry}
\usepackage{amsmath,amssymb}
\usepackage{enumitem}
\usepackage{hyperref}

\title{EM-Style Schema Refinement: Operations and Probability Calculations}
\author{}
\date{}

\begin{document}
\maketitle

EM is a classic algorithm that is typically used to jointly (1) update model parameters $\theta$ and (2) update latent variable assignments $z$ to data $x$ (e.g., EM for Gaussian mixtures). In this proposal, we run an EM-style loop where the latent variables are schema assignments and the ``parameters'' we update are not weights, but the schema itself.

On a high level, we do the following steps:
\begin{enumerate}[leftmargin=1.3em]
  \item \textbf{E-step:} assign labels to data based on the current (or initial) schema. Calculate metrics around how well these labels fit the data.
  \item \textbf{M-step:} update the schema---through split, merge, add, remove operations---then recompute metrics. Accept the new schema as the current schema if metrics improve.
\end{enumerate}
We iterate between the two steps until scores and assignments stabilize. We call this \emph{structural variational EM}.

We can also imagine the following extensions to this EM approach:
\begin{itemize}[leftmargin=1.3em]
  \item \textbf{Supervised Variant:} Instead of the E-step measuring how well the labels fit the datapoints (unsupervised variant), it can instead measure how well schema labels fit other labels (see Concept Bottleneck Models\footnote{\url{https://arxiv.org/pdf/2007.04612}}). This means we will tailor a schema towards a particular downstream variable.
  \item \textbf{Causal Variant:} We can extend the E-step to measure how well the schema captures causality by measuring, for example, interventional log-likelihood, invariance penalties $\mathrm{Inv}$, and yield identifiable, larger (C)ATEs. This means we will tailor a schema towards a particular causal variable.
\end{itemize}

\section{Setup}
Given observed texts $X = \{x_i\}_{i=1}^N$, an initial schema $Z_0$, and a reliable-enough mechanism $Z_0=\{z_j\}_{j=1}^K$ to apply this to all texts (i.e., our initial methodology).\footnote{Texts may be single- or multi-label.}

\section{E-Step (Labeling / Scoring)}

\subsection{Label Assignment and Probability Calculation}
We calculate several different probabilities to help us measure how well a label fits a text.
\begin{align}
\log p(z_i \mid x) &= \sum_{t \in \text{text tokens of } z_j} \log p(z_i \mid z_j, x), \tag{1} \\
\log p(x_i) &= \sum_{t \in \text{text tokens of } x_i} \log p(x_{i,t}), \tag{2} \\
p(x_i \mid z_j) &= \frac{p(z_j \mid x_i)\, p(z_j)}{\sum_{k=1}^K p(z_k \mid x_i)\, p(z_k)}. \tag{3}
\end{align}
Where $p(\cdot)$ are probabilities generated by an LLM. There are many ways to calculate the calibrated probability $p(z_j \mid x_i)$ (it has to be calibrated, because some labels $z$ are just more likely to be generated anyway). One approach is described in \url{https://arxiv.org/pdf/2102.09690} and implemented in \url{https://github.com/alex2awesome/schema-generation/blob/main/src/utils_probability_calibrator.py}.

We can use $p(z_j \mid x_i)$ to assign labels to texts. Then, we can use these probabilities to calculate statistics.

\subsection{Document-Level Scores}
Given probabilities for $x$, $z$, $x \mid z$, $z \mid x$, we compute scores across a document given a label scheme.

\paragraph{Average log-likelihood per category.}
\begin{equation}
L(z_j) = \frac{1}{|X_{z_j}|} \sum_{x_i \in X_{z_j}} \log p(x_i \mid z_j). \tag{4}
\end{equation}

\paragraph{Average posterior confidence.}
\begin{equation}
C(z_j) = \frac{1}{N} \sum_{i=1}^N p(z_j \mid x_i). \tag{5}
\end{equation}

\paragraph{Embedding centroid and intra-category variance.}
Let $e_{x_i}$ be a text embedding.
\begin{equation}
e_{z_j} = \frac{1}{|X_{z_j}|} \sum_{x_i \in X_{z_j}} e_{x_i},
\qquad
V(z_j) = \frac{1}{|X_{z_j}|} \sum_{x_i \in X_{z_j}} \|e_{x_i} - e_{z_j}\|^2. \tag{6}
\end{equation}

\paragraph{Pairwise category similarity.}
\begin{equation}
\mathrm{sim}(z_a, z_b) = \frac{e_{z_a} \cdot e_{z_b}}{\|e_{z_a}\|\, \|e_{z_b}\|}. \tag{7}
\end{equation}

\paragraph{Baseline.}
\begin{equation}
L_{\mathrm{baseline}} = \frac{1}{N} \sum_{i=1}^N \log p(x_i). \tag{8}
\end{equation}

\paragraph{Poorly explained texts.}
\begin{equation}
p_{\max}(x_i \mid Z) = \max_{1 \le j \le K} p(x_i \mid z_j). \tag{9}
\end{equation}

\subsection{Corpus-Level Scores}
We can also calculate corpus-level scores to understand the schema at timestep $t$, $V_t$, overall. Let $D_{\mathrm{train}}$ be the working set and $D_{\mathrm{test}}$ held out.

\paragraph{Total (conditional) log-likelihood.}
\begin{equation}
\log L_{\mathrm{cond}}(D) = \sum_{x_i \in D} \log p\!\left(x_i \mid \hat{z}_i\right). \tag{10}
\end{equation}

\paragraph{AIC \& BIC.}
Let $k$ be schema complexity (e.g., $k=K$ or a richer count) and $n = |D_{\mathrm{test}}|$:
\begin{align}
\mathrm{AIC} &= 2k - 2 \log L_{\mathrm{cond}}(D_{\mathrm{test}}), \tag{11} \\
\mathrm{BIC} &= k \log n - 2 \log L_{\mathrm{cond}}(D_{\mathrm{test}}). \tag{12}
\end{align}

\paragraph{Perplexity.}
Let $T = \sum_{x_i \in D} |x_i|$ be the total token count:
\begin{equation}
\mathrm{PPL}(D) = \exp\!\left(-\frac{1}{T} \sum_{x_i \in D} \log p\!\left(x_i \mid \hat{z}_i\right)\right). \tag{13}
\end{equation}

\paragraph{ELBO (variational framing).}
With $q_{ij} = p(z_j \mid x_i)$ as a variational posterior and prior $p(z_j)$:
\begin{equation}
\mathcal{L} = \sum_{i=1}^N \sum_{j=1}^K q_{ij} \log p(x_i, z_j) + H(Q),
\quad \text{where } p(x_i, z_j) = p(x_i \mid z_j)\, p(z_j). \tag{14}
\end{equation}

\section{M-Step (Schema Update Proposals)}
Given scores calculated during the E-step, we decide how to update our schema, yielding the current schema $Z_t$.

\subsection{Structural Moves}
\begin{itemize}[leftmargin=1.3em]
  \item \textbf{Split $z_j$.} When $L(z_j)$ is in the bottom quartile and/or $C(z_j)$ in the bottom quartile and/or $V(z_j)$ in the top quartile: cluster $X_{z_j}$ (e.g., $k=2$) into $X_{z_{j1}}, X_{z_{j2}}$; form $z_{j1}, z_{j2}$.
  \item \textbf{Merge $z_a, z_b$.} When $\mathrm{sim}(z_a, z_b) > \tau_{\mathrm{sim}}$ and $|L(z_a)-L(z_b)| < \varepsilon_L$, $|C(z_a)-C(z_b)| < \varepsilon_C$. Pool members into $z_{a\cup b}$, recompute quantities, and score.
  \item \textbf{Remove $z_j$.} When $|X_{z_j}|/|X| < \tau_{\min\text{-}size}$ or $L(z_j) \le L_{\mathrm{baseline}}(1+\delta)$. Drop $z_j$, reassign its texts by $p(z\mid x)$ over remaining labels, recompute, and score.
  \item \textbf{Add $z_{\mathrm{new}}$.} Cluster worst-explained texts (e.g., bottom 5--10\% by $p_{\max}(x_i \mid Z)$); introduce $z_{\mathrm{new}}$, recompute, and score.
\end{itemize}

\subsection{Scoring \& Acceptance}
Recalculate scores across the new schema $V_t$. Accept the proposed new schema $Z_t$ if $\mathrm{BIC}$ (or $\mathrm{AIC}$) decreases on $D_{\mathrm{test}}$ or if $\mathcal{L}$ increases. After accepting a schema change, re-run the E-step.

\section{Stopping Criteria}
Stop when any holds:
\begin{enumerate}[leftmargin=1.3em]
  \item \textbf{Score convergence:} $\left|\mathrm{Score}(t{+}1)-\mathrm{Score}(t)\right| / \left|\mathrm{Score}(t)\right| < \varepsilon$ (e.g., $\varepsilon = 0.01$).
  \item \textbf{Schema stability:} fewer than, say, 5\% of items change labels and no structural move is accepted.
  \item \textbf{Generalization plateaus:} held-out perplexity no longer improves.
  \item \textbf{Max iterations:} e.g., 5--10 passes.
\end{enumerate}

\end{document}
